\documentclass[11pt,a4paper]{article}
\usepackage{amsmath,amssymb,amsthm}
\usepackage[margin=2.5cm]{geometry}
\usepackage{hyperref}

\newtheorem{definition}{Definition}[section]
\newtheorem{theorem}[definition]{Theorem}
\newtheorem{proposition}[definition]{Proposition}
\newtheorem{remark}[definition]{Remark}
\newtheorem{conjecture}[definition]{Conjecture}

\title{Hyperseed Formalization:\\
  Provably Safe AGI Is Potentially Dangerous}
\author{ProtomegaTron (automated formalization)}
\date{2026-09-18}

\begin{document}
\maketitle

\begin{abstract}
We formalize Ben Goertzel's ``Provably Safe AGI Is Potentially Dangerous''
(2024-05-21) within the Hyperseed ontology. Key mappings: the specification
problem as fiber definition, the model-reality gap as fiber-base gap,
intractability as fiber complexity, false security as fiber-base confusion,
and defense in depth as multi-fiber verification.
\end{abstract}

\tableofcontents

\section{Fiber definition: the specification problem}

\begin{theorem}[Specification = fiber definition --- claims 1--4, 21]
You can't verify the fiber without first defining it. Specifying ``safe''
behavior for AGI requires formalizing human values --- complex, contextual,
contradictory, and evolving:
\[
  \text{Verify}(F) \text{ requires } \text{Define}(F)
  \quad \text{but} \quad \text{Define}(F_{\text{safe}}) = \text{unsolved}
\]
Even if specified today, the specification drifts as society evolves.
\end{theorem}

\section{Fiber-base gap: model vs reality}

\begin{theorem}[Model gap --- claims 9--12, 22, 24]
Proofs about formal models (fiber) don't guarantee properties of real
systems (base). The gap is where failures hide:
\[
  \text{Prove}(F_{\text{model}} \models \phi) \not\Rightarrow
  B_{\text{real}} \models \phi
\]
The false sense of security comes from confusing fiber verification
(proof holds in model) with base verification (system is actually safe).
This is a category error with potentially catastrophic consequences.
\end{theorem}

\section{Fiber complexity: intractability}

\begin{proposition}[Complexity barrier --- claims 5--8, 23]
Formal verification becomes intractable when the fiber is too complex.
Self-modification makes the fiber dynamic:
\[
  \text{Verify}(F_t) \quad \text{but} \quad F_{t+1} = \text{modify}(F_t)
  \implies \text{Verify}(F_{t+1}) \text{ needed again}
\]
The modification space is unbounded, making exhaustive verification
impossible. Any tractable verification uses approximations that
introduce gaps.
\end{proposition}

\section{Multi-fiber verification: defense in depth}

\begin{proposition}[Complementary approach --- claims 17--19, 25]
Multiple independent safety mechanisms provide multi-fiber verification:
\[
  \text{Safe} \approx \bigcap_i \text{Verified}_i(F_i)
\]
No single mechanism is sufficient; the intersection of partial
verifications is more reliable than any single verification alone.
\end{proposition}

\section{Cross-references}

\begin{remark}[Connections]
\begin{itemize}
\item \textbf{Seven Flavors of AGI Catastrophe} (2026-06-23): systematic
  analysis of failure modes that ``provable safety'' claims to prevent.
\item \textbf{Why Everyone Dies Gets AGI All Wrong} (2025-10-01): defense
  in depth through hybrid architecture.
\item \textbf{Evolving Deeply Ethical AGI} (2026-01-06): the specification
  problem connects to non-dual motivation --- safety can't be specified
  as a single objective.
\end{itemize}
\end{remark}

\end{document}
